\clearpage
\phantomsection
\vspace*{1cm}
\begin{center}
    \textbf{\fontsize{20}{24}\selectfont Abstract}
\end{center}
\vspace{1cm}
% Nội dung tóm tắt

Amid the digital transformation of the finance and banking sector and the growth of online lending, building credit risk prediction systems that offer high accuracy, transparent interpretability, and practical decision-support value has become an urgent requirement. However, while most modern machine learning models achieve strong predictive performance, they remain limited in explaining the causes leading to default risk, which reduces their trustworthiness and applicability in financial risk management settings.

This study proposes a credit risk prediction framework based on causal machine learning combined with Bayesian networks, implemented in a logical and rigorous manner through five advanced integrated stages: balancing the data with the SMOTE algorithm to optimise sample quality (S\_Pro), screening and selecting strongly discriminative features via Lasso regression (L\_Pro), discretising continuous variables with the K-means algorithm (K\_Pro), learning the network structure and parameters through the Tabu Search algorithm combined with Maximum Likelihood Estimation (MLE) (BN\_Pro), and finally performing forward/backward inference together with sensitivity analysis to fully exploit the causal structure (In\_Pro). In addition, the study integrates model explanation techniques such as SHAP, LIME, sensitivity analysis, forward–backward reasoning, classification threshold optimisation, and Decision Curve Analysis in order to comprehensively evaluate the model's applied effectiveness.

The experimental results show that the model achieves superior discrimination of high-risk customers compared with many traditional methods, while also providing clear causal relationships between financial factors and the likelihood of default. Optimising the decision threshold helps to effectively balance profit against the cost of risk, whereas Decision Curve Analysis demonstrates the model's practical economic value across a range of credit-granting scenarios. The study contributes a transparent, interpretable and decision-supporting credit risk assessment framework that helps financial institutions improve appraisal quality, reduce credit losses, and optimise risk management effectiveness in the modern business environment.
\vspace*{\fill} % Lò xo đẩy lên (cân bằng với lò xo ở trên)